\documentclass{article}

\usepackage{booktabs}
\usepackage{amsfonts,amsmath,graphicx} %subfigure,
\usepackage{blindtext}
\usepackage{outlines}
\usepackage{comment}
\usepackage{amsmath}
\usepackage{hyperref}
\usepackage[dvipsnames]{xcolor}



\title{Interactive Visual Exploration of ODE Parameter Spaces through VAE Embeddings}
\author{Emiliano Islas and Noel Martin Poothokaran}

\begin{document}

\maketitle


\section{Problem Description \& Motivation}


Dynamical systems are often modeled by ordinary differential equations whose
form is abstractly represented as
\begin{equation}
    \dot{x} = f(x;\,\theta),
\end{equation}
where $\theta$ denotes a set of model parameters. In practice, the influence
of individual parameters on the qualitative behavior of the solution is
difficult to interpret, particularly when parameters interact nonlinearly.
Variational Autoencoders (VAEs) can encode high-dimensional data into a
low-dimensional continuous latent space, and have been shown to discover
interpretable structure in an unsupervised manner. 

The primary contribution of this work is a visual analysis framework
for parameterized dynamical systems. Rather than examining individual
time-series plots or conducting parameter sweeps by repeatedly solving the ODE,
the trained VAE provides an interactive low-dimensional map in which the user
can visually identify how trajectories cluster, separate, and transition as
parameters change. This approach transforms parameter exploration from a
computational task---solve, plot, compare---into a navigable visual space where
the relationship between parameters and dynamics can be perceived at a glance. Our tool also aims to generalize to high-dimensional situations such as a large number of parameters $\theta$ while retaining visual interpretability.


\section{Related Work}

Our work draws on three lines of research: latent variable models for
dynamical data, visual parameter space analysis, and interactive exploration of
learned embeddings.

\paragraph{Variational Autoencoders.}
Kingma and Welling~\cite{kingma2014auto} introduced the VAE as a generative
model that couples an encoder--decoder architecture with a principled
variational inference objective. The continuous latent space and the existence
of a decoder distinguish the VAE from non-parametric projections such as
t-SNE~\cite{vandermaaten2008visualizing}, making it possible to both visualize
and generate data. We adopt the VAE as the core representation learning
component of our framework.

\paragraph{Neural ODEs and Latent Dynamical Models.}
Chen et al.~\cite{chen2018neural} showed that continuous-depth residual
networks can be interpreted as ODEs, establishing a formal bridge between deep
learning and dynamical systems. Rubanova et
al.~\cite{rubanova2019latent} extended this idea by combining a VAE with an
ODE-based latent dynamics model for irregularly sampled time series. While
their goal is prediction and interpolation, our objective is interactive visual
exploration of the mapping from parameters to qualitative dynamical behavior.

\paragraph{Visual Parameter Space Analysis.}
Sedlmair et al.~\cite{sedlmair2014visual} proposed a conceptual framework for
visual parameter space analysis, systematizing how analysts explore the
relationship between input parameters and model outputs through visualization.
Our work instantiates this framework using learned latent representations
rather than direct sampling grids or simple projection methods, enabling
scalable exploration of high-dimensional parameter spaces.

\paragraph{Interactive Exploration of Learned Embeddings.}
Liu et al.~\cite{liu2018visual} demonstrated interactive visual exploration of
neural word embeddings, linking latent-space structure to semantic
relationships. Their coordinated-view design informs our interface, though we
operate in a fundamentally different domain---ODE trajectories rather than word
vectors---and additionally exploit the VAE decoder for generative
interpolation.

\paragraph{Dimensionality Reduction for Visualization.}
t-SNE~\cite{vandermaaten2008visualizing} and UMAP~\cite{mcinnes2018umap} are
widely used to project high-dimensional data to two dimensions for visual
inspection. We include these methods as baselines; however, because they lack a
decoder, they cannot support the generative trajectory exploration that is a
key feature of our system.

\paragraph{Interactive Latent Space Visualization Tools.}
Taylor Denouden's VAE Latent Space
Explorer~\cite{denouden2017vaexplorer} provides a compelling demonstration of
interactive latent space navigation for image data: users click within a
two-dimensional latent grid and observe decoded outputs in real time, making
the structure of the learned manifold directly perceptible. Our generative
interpolation panel draws inspiration from this interface, adapting the
click-to-decode paradigm from image generation to the visualization of ODE
trajectories. Whereas the VAE Latent Space Explorer displays reconstructed
images on a static grid, our system embeds the interaction within a
coordinated multiple-view environment and links decoded trajectories to
parameter-space and time-series views, enabling analytical reasoning about
dynamical behavior rather than passive inspection.
    

\section{Planned Technical Approach}

Our framework consists of four stages: data generation, representation
learning, interactive visualization, and scalability support.

\subsection{Data Generation}

We select canonical parameterized ODE systems that exhibit rich qualitative
transitions as parameters vary. Initial benchmark systems include the Lorenz
system (fixed points, limit cycles, and chaotic attractors), the
FitzHugh--Nagumo model (excitable and oscillatory regimes), and the damped
Duffing oscillator (periodic and chaotic responses). For each system, we sample
parameter vectors $\theta$ from the feasible region using Latin hypercube
sampling to ensure uniform coverage. Each sampled $\theta$ is numerically
integrated using a standard solver (e.g., Runge--Kutta~4/5) to produce a
trajectory. Trajectories are then represented as fixed-length feature vectors
by uniform temporal resampling, optionally augmented with summary statistics
such as dominant spectral frequencies or estimated Lyapunov exponents.

\subsection{VAE Architecture and Training}

A Variational Autoencoder is trained to map trajectory feature vectors to a
low-dimensional latent space $\mathbf{z} \in \mathbb{R}^d$, with $d = 2$ or
$d = 3$ for direct visual mapping. The encoder
$q_\phi(\mathbf{z}\mid\mathbf{x})$ and decoder
$p_\psi(\mathbf{x}\mid\mathbf{z})$ are parameterized as fully connected
networks. Training minimizes the standard evidence lower bound (ELBO) loss:
\begin{equation}
    \mathcal{L}(\phi,\psi) \;=\;
    -\,\mathbb{E}_{q_\phi(\mathbf{z}\mid\mathbf{x})}
        \!\bigl[\log p_\psi(\mathbf{x}\mid\mathbf{z})\bigr]
    \;+\;
    \mathrm{KL}\!\bigl(q_\phi(\mathbf{z}\mid\mathbf{x})
        \,\|\, p(\mathbf{z})\bigr).
\end{equation}
To encourage disentangled latent dimensions that align with interpretable
dynamical properties, we experiment with the $\beta$-VAE
variant~\cite{higgins2017betavae}, which up-weights the KL term. We also
explore conditioning the decoder on $\theta$ so that the latent space captures
variation beyond what the parameters already explain.

\subsection{Interactive Visual Analytics Interface}

We implement a web-based coordinated multiple-view system using D3.js and
React. The interface comprises the following linked views:

\begin{itemize}
    \item \textbf{Latent Space View.} A scatterplot of encoded trajectories in
    the two- or three-dimensional latent space. Points are colored by parameter
    values, cluster membership, or user-defined attributes. Brushing in this
    view propagates selections to all other views.

    \item \textbf{Trajectory Detail View.} Displays time-series plots or phase
    portraits corresponding to selected latent-space points, enabling direct
    inspection of the dynamical behavior underlying each region.

    \item \textbf{Parameter Space View.} Visualizes the parameter coordinates
    of selected trajectories using parallel coordinates or a scatterplot
    matrix, revealing which parameter regimes map to which latent regions.

    \item \textbf{Generative Interpolation Panel.} The user clicks or drags a
    cursor through the latent space; the decoder generates reconstructed
    trajectories in real time, revealing smooth transitions between dynamical
    regimes and surfacing behaviors not present in the training sample.
    By colour-coding the latent space by known parameter values, an analyst can
    immediately see which directions in the representation correspond to which
    physical quantities. Sweeping along these directions and decoding the
    result produces a continuous animation of the evolving waveform, offering
    an intuitive understanding of parameter sensitivity that static plots
    cannot provide.

    \item \textbf{Cluster and Boundary Overlay.} Automatic clustering
    (DBSCAN or Gaussian mixture models) in the latent space delineates
    qualitative behavioral regimes. Decision boundaries are rendered as
    contours, highlighting regions that may correspond to bifurcations in the
    original system.
\end{itemize}

All views are linked via brushing and filtering: selecting a cluster in the
latent space highlights the corresponding parameter ranges and trajectory
shapes, and adjusting a parameter slider animates the movement of a point
through the latent space.

\subsection{Scalability to High-Dimensional Parameter Spaces}

Because the VAE encoder maps trajectories to $\mathbf{z}$ independently of the
dimensionality of $\theta$, the latent space view remains two- or
three-dimensional even when the ODE has many parameters. To maintain
interpretability of $\theta$ itself, we supplement the interface with a
parallel-coordinates view and support dimensionality reduction of the parameter
space (e.g., PCA or UMAP) as an auxiliary panel.

\section{Dataset(s)}

Because our framework targets ODE models rather than fixed empirical datasets,
the data are generated synthetically by solving the governing equations over
sampled parameter configurations. This approach provides full control over
ground-truth labels (e.g., known bifurcation boundaries) and allows scaling to
arbitrary sample sizes.

\begin{itemize}
    \item \textbf{Lorenz system} ($\theta = [\sigma, \rho, \beta]$, 3
    parameters): a canonical testbed for chaotic dynamics, with well-documented
    transitions between fixed points, periodic orbits, and strange attractors.

    \item \textbf{FitzHugh--Nagumo model} ($\theta = [a, b, \tau, I]$, 4
    parameters): a simplified model of neuronal excitability, exhibiting
    excitable, oscillatory, and bistable regimes.

    \item \textbf{Lotka--Volterra predator--prey model}
    ($\theta = [\alpha, \beta, \delta, \gamma]$, 4 parameters): governed by
    \begin{align}
        \frac{dx_1}{dt} &= \alpha\, x_1 - \beta\, x_1 x_2, \nonumber\\
        \frac{dx_2}{dt} &= \delta\, x_1 x_2 - \gamma\, x_2,
    \end{align}
    this system produces periodic oscillations whose shape, amplitude, and
    frequency depend on the parameters and initial conditions. It serves as an
    accessible and widely studied benchmark whose closed orbits provide clear
    visual signatures for the latent space to capture.

    \item \textbf{Damped Duffing oscillator} ($\theta = [\delta, \alpha,
    \beta, \gamma, \omega]$, 5 parameters): demonstrates period-doubling
    routes to chaos, providing a higher-dimensional parameter space to test
    scalability.
\end{itemize}

For each system, we generate on the order of $10{,}000$ trajectories by
sampling parameters via Latin hypercube sampling. All data generation code will
be released alongside the tool.

\section{Evaluation Methodology}

\subsection{Expected Results}

We expect the trained VAE to produce a latent space in which trajectories with
similar qualitative behavior (e.g., periodic, chaotic) form coherent clusters,
and in which transitions between regimes appear as smooth gradients or sharp
boundaries that correspond to known bifurcations. The interactive interface
should enable a user to (1)~identify qualitative behavioral regimes without
manual trajectory inspection, (2)~locate parameter-space boundaries associated
with bifurcations, and (3)~generate and inspect novel trajectories by
navigating the latent space, thereby accelerating hypothesis formation about
parameter--dynamics relationships.

\subsection{Evaluation Methodology}

We evaluate the framework along three complementary axes.

\paragraph{Quantitative Evaluation.}
We measure VAE reconstruction error on held-out trajectories and assess latent
space quality using silhouette scores computed against ground-truth behavioral
labels (e.g., fixed point vs.\ limit cycle vs.\ chaos). We compare the VAE
embedding against t-SNE and UMAP baselines in terms of cluster separation and
neighborhood preservation.

\paragraph{Case Studies.}
We conduct two detailed case studies---one on the Lorenz system and one on the
Lotka-Volterra model---demonstrating the end-to-end workflow: a user explores
the latent space, identifies clusters, examines decoded trajectories, and
relates findings to known bifurcation diagrams from the dynamical-systems
literature.

\paragraph{Qualitative User Feedback.}
Time permitting, we solicit informal feedback from graduate students with
backgrounds in dynamical systems or machine learning, focusing on the
interpretability of the latent space, the utility of generative interpolation,
and the overall usability of the coordinated-view interface.


% ============================================================
% BIBLIOGRAPHY (BibTeX entries)
% ============================================================
% Place the following in your .bib file.

\begin{thebibliography}{9}

\bibitem{kingma2014auto}
D.~P. Kingma and M.~Welling,
``Auto-Encoding Variational Bayes,''
in \emph{Proceedings of the 2nd International Conference on Learning
Representations (ICLR)}, 2014.
Available: \url{https://arxiv.org/abs/1312.6114}

\bibitem{vandermaaten2008visualizing}
L.~van~der~Maaten and G.~Hinton,
``Visualizing Data using t-SNE,''
\emph{Journal of Machine Learning Research}, vol.~9, pp.~2579--2605, 2008.
Available: \url{https://jmlr.org/papers/v9/vandermaaten08a.html}

\bibitem{chen2018neural}
R.~T.~Q. Chen, Y.~Rubanova, J.~Bettencourt, and D.~Duvenaud,
``Neural Ordinary Differential Equations,''
in \emph{Advances in Neural Information Processing Systems (NeurIPS)}, vol.~31,
2018.
Available: \url{https://arxiv.org/abs/1806.07366}

\bibitem{rubanova2019latent}
Y.~Rubanova, R.~T.~Q. Chen, and D.~Duvenaud,
``Latent Ordinary Differential Equations for Irregularly-Sampled Time Series,''
in \emph{Advances in Neural Information Processing Systems (NeurIPS)}, vol.~32,
2019.
Available: \url{https://arxiv.org/abs/1907.03907}

\bibitem{sedlmair2014visual}
M.~Sedlmair, C.~Heinzl, S.~Bruckner, H.~Piringer, and T.~M{\"o}ller,
``Visual Parameter Space Analysis: A Conceptual Framework,''
\emph{IEEE Transactions on Visualization and Computer Graphics}, vol.~20,
no.~12, pp.~2161--2170, 2014.
Available: \url{https://doi.org/10.1109/TVCG.2014.2346321}

\bibitem{liu2018visual}
S.~Liu, P.-T. Bremer, J.~J. Thiagarajan, V.~Srikumar, B.~Wang, Y.~Livnat,
and V.~Pascucci,
``Visual Exploration of Semantic Relationships in Neural Word Embeddings,''
\emph{IEEE Transactions on Visualization and Computer Graphics}, vol.~24,
no.~1, pp.~553--562, 2018.
Available: \url{https://doi.org/10.1109/TVCG.2017.2745141}

\bibitem{mcinnes2018umap}
L.~McInnes, J.~Healy, and J.~Melville,
``UMAP: Uniform Manifold Approximation and Projection for Dimension
Reduction,''
\emph{arXiv preprint arXiv:1802.03426}, 2018.
Available: \url{https://arxiv.org/abs/1802.03426}

\bibitem{higgins2017betavae}
I.~Higgins, L.~Matthey, A.~Pal, C.~Burgess, X.~Glorot, M.~Botvinick,
S.~Mohamed, and A.~Lerchner,
``$\beta$-VAE: Learning Basic Visual Concepts with a Constrained Variational
Framework,''
in \emph{Proceedings of the 5th International Conference on Learning
Representations (ICLR)}, 2017.
Available: \url{https://openreview.net/forum?id=Sy2fzU9gl}

\bibitem{denouden2017vaexplorer}
T.~Denouden,
``VAE Latent Space Explorer,''
2017. [Online]. Available: \url{https://tayden.github.io/VAE-Latent-Space-Explorer/}

\end{thebibliography}

\end{document}